\documentclass[../main.tex]{subfiles}
\begin{document}
\chapter{2011-2012学年微积分（一）（上）期中考试}

\section{基本计算题(每小题 5 分，共 60 分)}
\begin{enumerate}
    \item 计算极限$l=\lim_{n\to\infty}\frac{2011^{n}}{n!}$.
          \vspace{7em}

    \item 设$x\to0$时，无穷小量$u=\sqrt[4]{1-a\arctan x^{2}}-1$与$v=\ln\cos x$等价，求常数$a$的值.
          \vspace{7em}

    \item 计算极限$l=\lim_{x\to 0}\frac{\mathrm{e}^{x}-\mathrm{e}^{\sin x}}{x^{3}}$.
          \vspace{8em}

    \item 计算极限$l=\lim_{x\to 0}\left(\frac{(1+x)^{\frac{1}{x}}}{\mathrm{e}}\right
              )^{\frac{1}{x}}$.
          \vspace{8em}

    \item 设曲线$y=f(x)$与$y=\arctan 2x$在原点相切，求$\lim_{n\to\infty}nf\left(\frac{1}{n}
              \right)$.
          \vspace{9em}

    \item 设函数$\varphi(x)$在$x=1$处连续，且任给自然数$n$，有$\varphi\left(\frac{n}{n+1}
              \right)=\sqrt[n]{n}$.

          （1）求$\varphi(1)$；

          （2）设$f(x)=(x-1)\varphi(x)$，求$f'(1)$.
          \vspace{9em}

    \item 设$y=\sqrt{x+\sqrt{x}}$，求$\mathrm{d}y(1)$.
          \vspace{9em}

    \item 设函数$y=y(x)$由方程$x^{2}-xy+y^{2}=1$确定，求$y'$，$y''$.
          \vspace{9em}

    \item 设$\begin{cases}
                  x=\ln\cos t, \\
                  y=\sin t-t\cos t
              \end{cases}$，求$\frac{\mathrm{d}y}{\mathrm{d}x}$和$\left.\frac{\mathrm{d}^{2}
                  y}{\mathrm{d}x^{2}}\right|_{t=\frac{\pi}{4}}$.
          \vspace{9em}

    \item 设$f(x)=\frac{x+1}{2x^{2}+5x+2}$，求$f^{(n)}(x)(n>1)$.
          \vspace{8em}

    \item 确定自然数$n$的范围，使$f(x)=
              \begin{cases}
                  x^{n} \sin\frac{1}{x}, & x\neq0, \\
                  0,                     & x=0
              \end{cases}$的导函数$f'(x)$在$x=0$处连续.
          \vspace{8em}

    \item 设函数$f(x)=\mathrm{e}^{-\frac{1}{x^{2}}}\arctan\frac{x+1}{x-1}$，指出其间断点，并说明间断点的类型.
          \vspace{8em}
\end{enumerate}

\section{综合题(每小题 6 分，共 24 分)}
\begin{enumerate}
    \setcounter{enumi}{12}
    \item 设$f(x)$在$[0,2]$上连续，在$(0,2)$内可微，且$f(0)\cdot f(2)>0$，$f
              (0)\cdot f(1)<0$，证明：存在$\xi\in(0,2)$，使$f'(\xi)=f(\xi)$.
          \vspace{9em}

    \item 设函数$y=f(x)$和它的反函数$x=\varphi(y)$均存在三阶导数，且$y'\neq0$，请推导出反函数的求导公式$\frac{\mathrm{d}x}{\mathrm{d}y}$，$\frac{\mathrm{d}^{2}
                  x}{\mathrm{d}y^{2}}$，$\frac{\mathrm{d}^{3} x}{\mathrm{d}y^{3}}$【用$y'$，$y
              ''$，$y'''$表示】.
          \vspace{9em}

    \item 证明：当$x\geq1$时，有$\arctan x=\frac{\pi}{4}+\frac{1}{2}\arccos\frac{2x}{1+x^{2}}$.
          \vspace{10em}

    \item 设函数$f(x)$在$[a,+\infty)$上连续，在$(a,+\infty)$内可导，且$f'(x)
              >1$.若$f(a)<0$，证明：方程$f(x)=0$在区间$(a,a-f(a))$内有唯一根.
          \vspace{11em}
\end{enumerate}

\section{证明题(每小题 8 分，共 16 分)}
\begin{enumerate}
    \setcounter{enumi}{16}
    \item 分别叙述数列$x_{n}$有界和收敛（以$\lim_{n\to\infty}x_{n}=a$为例）的定义，并证明：收敛数列是有界数列.
          \vspace{12em}

    \item 如果记$\xi=\theta\cdot x$，$0<\theta<1$，则Lagrange中值公式$f(x)-f
              (0)=xf'(\xi)$可以写作：$f(x)-f(0)=xf'(\theta x)$，$0<\theta<1$，$\theta$的大小通常与$x$相关.

          （1）若$f''(0)\neq0$，试证：$\lim_{x\to0}\theta=\frac{1}{2}$；

          （2）设$f(x)=\arctan x$，求$\lim_{x\to0}\theta$.
          \vspace{12em}
\end{enumerate}

\chapter{2011-2012学年微积分（一）（上）期中考试参考答案}
\section{基本计算题(每小题 5 分，共 60 分)}
\begin{enumerate}
    \item \textbf{Solution}. 当$n>2011$时，$\frac{a_{n}}{a_{n-1}}=\frac{2011^{n}}{n!}
              \cdot\frac{(n-1)!}{2011^{n-1}}=\frac{2011}{n}<1$，

          数列$\{a_{n}\}$从第$2012$项开始单调递减，且$a_{n}>0$，所以数列$\{a_{n}\}$收敛.

          对$a_{n} = \frac{2011}{n}a_{n-1}$两边取极限，得$l=0$.

    \item \textbf{Solution}. 由题意可知
          \[
              \begin{aligned}
                  1 = {} & \lim_{x\to 0}\frac{\sqrt[4]{1-a\arctan x^2}-1}{\ln\cos x}                                                  \\
                  = {}   & \lim_{x\to0}\frac{-\frac{1}{4}a\arctan x^2}{\cos x-1}=\lim_{x\to 0}\frac{-\frac{a}{4} x^2}{-\frac{x^2}{2}} \\
                  = {}   & \frac{a}{2}.
              \end{aligned}
          \]
          故$a=2$.

    \item \textbf{Solution}. 由Lagrange中值定理，存在$\xi$介于$x$和$\sin x$之间，使得$\mathrm{e}
              ^{x}-\mathrm{e}^{\sin x}=\mathrm{e}^{\xi}(x-\sin x)$.

          易知当$x\to0$时，$\xi\to0$，所以
          \[
              \begin{aligned}
                  l= {} & \lim_{x\to0}\mathrm{e}^{\xi}\frac{x-\sin x}{x^3} \\
                  = {}  & \lim_{x\to0}\frac{1-\cos x}{3x^2}                \\
                  = {}  & \frac{1}{6}.
              \end{aligned}
          \]

    \item \textbf{Solution}.
          \[
              \begin{aligned}
                  l = {} & \exp\left[\lim_{x\to0}\frac{\ln\frac{(1+x)^{\frac{1}{x}}}{\mathrm{e}}}{x}\right]=\exp\left[\lim_{x\to0}\frac{\frac{1}{x}\ln(1+x)-1}{x}\right] \\
                  = {}   & \exp\left[\lim_{x\to0}\frac{\ln(1+x)-x}{x^2}\right]=\exp\left[\lim_{x\to0}\frac{x-\frac{x^2}{2}+o(x^2)-x}{x^2}\right]                         \\
                  = {}   & \frac{1}{\sqrt{\mathrm{e}}}.
              \end{aligned}
          \]

    \item \textbf{Solution}. 由题意可知
          \[
              f(0)=y(0)=0,\quad f'(0)=y'(0)=\left.\frac{2}{1+4x^{2}}\right|_{x=0}=2,
          \]
          所以
          \[
              \lim_{n\to\infty}nf\left(\frac{1}{n}\right)=\lim_{n\to\infty}\frac{f\left(\frac{1}{n}\right)-f(0)}{\frac{1}{n}-0}=f'(0)=2.
          \]

    \item \textbf{Solution}. （1）因$\varphi(x)$在$x=1$处连续，所以
          \[
              \varphi(1)=\lim_{n\to\infty}\varphi\left(\frac{n}{n+1}\right)=\lim_{n\to\infty}\sqrt[n]{n}=1.
          \]

          （2）由导数的定义可得，
          \[
              f'(1)=\lim_{x\to1}\frac{f(x)-f(1)}{x-1}=\lim_{x\to1}\frac{(x-1)\varphi(x)}{x-1}
              =\lim_{x\to1}\varphi(x)=\varphi(1)=1.
          \]

    \item \textbf{Solution}. 因为$y'=\frac{1}{2\sqrt{x+\sqrt{x}}}\cdot\left(1+\frac{1}{2\sqrt{x}}
              \right)$，所以$\mathrm{d}y(1)=y'(1)\,\mathrm{d}x=\frac{3\sqrt{2}}{8}\,\mathrm{d}
              x$.

    \item \textbf{Solution}. 方程两边对$x$求导，得$2x-y-x\cdot y'+2y\cdot y'=0$，所以$y
              '=\frac{2x-y}{x-2y}$.

          $y''=\frac{(2-y')(x-2y)-(2x-y)(1-2y')}{(x-2y)^{2}}=\frac{6}{(x-2y)^{3}}$.

    \item \textbf{Solution}. 因为
          \[
              \frac{\mathrm{d}y}{\mathrm{d}x}=\frac{\frac{\mathrm{d}y}{\mathrm{d}t}}{\frac{\mathrm{d}x}{\mathrm{d}t}}
              =\frac{\cos t+t\sin t-\cos t}{\frac{-\sin t}{\cos t}}=-t\cos t,
          \]
          \[
              \frac{\mathrm{d}^{2}y}{\mathrm{d}x^{2}}=\frac{\mathrm{d}}{\mathrm{d}t}\left
              (\frac{\mathrm{d}y}{\mathrm{d}x}\right)\cdot\frac{\mathrm{d}t}{\mathrm{d}x}
              =\left(-t\cos t\right)'\cdot\frac{1}{\frac{-\sin t}{\cos t}}=\frac{\cos t- t\sin
                  t}{\tan t},
          \]
          所以
          \[
              \left.\frac{\mathrm{d}^{2} y}{\mathrm{d}x^{2}}\right|_{t=\frac{\pi}{4}}
              =\frac{\sqrt{2}}{2}\left(1-\frac{\pi}{4}\right).
          \]

    \item \textbf{Solution}. 因为$f(x)=\frac{1}{3}\left(\frac{1}{2x+1}+\frac{1}{x+2}\right
              )$，所以
          \[
              f^{(n)}(x)=\frac{1}{3}\left[\frac{(-1)^{n}\,n!\,2^{n}}{(2x+1)^{n+1}}+\frac{(-1)^{n}\,n!}{(x+2)^{n+1}}
                  \right].
          \]

    \item \textbf{Solution}. 欲使$\lim_{x\to0}\frac{f(x)-f(0)}{x-0}=\lim_{x\to0}x
                  ^{n-1}\sin\frac{1}{x}$存在，必须$n-1>0$，此时$f'(0)=0$.

          欲使导函数$f'(x)$在$x=0$处连续，应有$\lim_{x\to0}f'(x)=f'(0)=0$.

          当$x\neq0$时，$f'(x)=nx^{n-1}\sin\frac{1}{x}-x^{n-2}\cos\frac{1}{x}$，所以必须$n
              -2>0$.

          综上所述，自然数$n$的范围是$n>2$.

    \item \textbf{Solution}. $f(x)$的间断点为$x=0$和$x=1$.

          因为
          \[\lim_{x\to0}f(x)=\lim_{x\to0}\mathrm{e}^{-\frac{1}{x^{2}}}\arctan\frac{x+1}{x-1}
              =\lim_{x\to0}\frac{\arctan\frac{x+1}{x-1}}{\mathrm{e}^{\frac{1}{x^2}}}=0,
          \]
          且$f(0)$无定义，所以$x=0$是$f(x)$的可去间断点；

          因为
          \[
              \begin{gathered}
                  \lim_{x\to1^-}f(x)=\lim_{x\to1^-}\mathrm{e}^{-\frac{1}{x^{2}}}\arctan\frac{x+1}{x-1}
                  =-\frac{\pi}{2\mathrm{e}},\\
                  \lim_{x\to1^+}f(x)=\lim_{x\to1^+}\mathrm{e}^{-\frac{1}{x^{2}}}
                  \arctan\frac{x+1}{x-1}=\frac{\pi}{2\mathrm{e}},
              \end{gathered}
          \]
          所以$x=1$是$f(x)$的跳跃间断点.
\end{enumerate}

\section{综合题(每小题 6 分，共 24 分)}
\begin{enumerate}
    \setcounter{enumi}{12}
    \item \textbf{Solution}. 由连续函数的介值定理，存在$\eta_{1}\in(0,1),\eta
              _{2}\in(1,2)$使得$f(\eta_{1})=f(\eta_{2})=0$.

          设$g(x)=\mathrm{e}^{-x}f(x)$，则$g(x)$在$[\eta_{1},\eta_{2}]$内可导，且$g(\eta
              _{1})=g(\eta_{2})=0$.

          由Rolle定理，存在$\xi\in(\eta_{1},\eta_{2})\subset(0,2)$使得$g'(\xi)=\mathrm{e}
              ^{-\xi}(f'(\xi)-f(\xi))=0$，所以$f'(\xi)=f(\xi)$.

    \item \textbf{Solution}.
          \[
              \frac{\mathrm{d}x}{\mathrm{d}y}=\frac{1}{\frac{\mathrm{d}y}{\mathrm{d}x}}
              =\frac{1}{y'},
          \]
          \[
              \frac{\mathrm{d}^{2} x}{\mathrm{d}y^{2}}=\frac{\mathrm{d}}{\mathrm{d}y}\left
              (\frac{1}{y'}\right)=\frac{\mathrm{d}}{\mathrm{d}x}\left(\frac{1}{y'}\right
              )\cdot\frac{\mathrm{d}x}{\mathrm{d}y}=\frac{-y''}{(y')^{2}}\cdot\frac{1}{y'}
              =-\frac{y''}{(y')^{3}},
          \]
          \[
              \frac{\mathrm{d}^{3} x}{\mathrm{d}y^{3}}=\frac{\mathrm{d}}{\mathrm{d}y}\left
              (-\frac{y''}{(y')^{3}}\right)=\frac{\mathrm{d}}{\mathrm{d}x}\left(-\frac{y''}{(y')^{3}}
              \right)\cdot\frac{\mathrm{d}x}{\mathrm{d}y}=-\frac{y'''(y')^{3}-3(y')^{2} (y'')^{2}}{(y')^{6}}
              \cdot\frac{1}{y'}=-\frac{y'y'''-3(y'')^{2}}{(y')^{5}}.
          \]

    \item \textbf{Solution}. 令$f(x)=\arctan x-\frac{1}{2}\arccos\frac{2x}{1+x^{2}}
              (x\geq1)$，则
          \[
              f'(x)=\frac{1}{1+x^{2}}+\frac{1}{2}\cdot\frac{\frac{2(1+x^{2})-4x^{2}}{(1+x^{2})^{2}}}{\sqrt{1-\left(\frac{2x}{1+x^2}\right)^{2}}}
              \equiv0,
          \]
          所以$f(x)$在$[1,+\infty)$上恒为常数.又$f(1)=\arctan 1-\frac{1}{2}\arccos 1=\frac{\pi}{4}$，所以当$x\geq1$时，
          \[
              \arctan x=\frac{\pi}{4}+\frac{1}{2}\arccos\frac{2x}{1+x^{2}}.
          \]

    \item \textbf{Solution}. 记$b=a-f(a)>a$，则由Lagrange中值定理，存在$\xi\in
              (a,b)$使得
          \[
              f(b)=f(a)+f'(\xi)(b-a)>f(a)+(b-a)=0.
          \]
          又由连续函数的介值定理，存在$\eta\in(a,b)$使得$f(\eta)=0$.

          因为$f(x)$是严格单调递增的，所以$\eta$是方程$f(x)=0$在区间$(a,a-f(a))$内的唯一根.
\end{enumerate}

\section{证明题(每小题 8 分，共 16 分)}
\begin{enumerate}
    \setcounter{enumi}{16}
    \item 数列$x_{n}$有界的定义：$\exists\,M>0$，$\forall\,n\in\mathbf{N}$，恒有$|
              x_{n}|\leq M$.

          数列$x_{n}$收敛于$a$的定义：$\forall\,\varepsilon>0$，$\exists\,N\in\mathbf{N}$，$\forall
              \,n>N$，恒有$|x_{n}-a|<\varepsilon$.

          \textbf{Proof}. 取$\varepsilon=1$，则存在自然数$N$使得当$n>N$时，恒有$|x_{n}
              -a|<1$，从而
          \[
              |x_{n}|=|x_{n}-a+a|<1+|a|.
          \]
          设$M=\max\{|x_{1}|,|x_{2}|,\cdots,|x_{N}|,1+|a|\}$，则$\forall\,n\in\mathbf{N}$，恒有$|
              x_{n}|\leq M$.所以数列$x_{n}$有界.

    \item \textbf{Proof}. （1）计算
          \[
              \begin{aligned}
                  f''(0)= {} & \lim_{x\to0}\frac{f'(x)-f'(0)}{x-0}= \lim_{x\to0}\frac{f'(\theta x)-f'(0)}{\theta x}                    \\
                  = {}       & \lim_{x\to0}\frac{\frac{f(x)-f(0)}{x}-f'(0)}{\theta x}= \lim_{x\to0}\frac{f(x)-f(0)-f'(0)x}{\theta x^2} \\
                  = {}       & \lim_{x\to0}\frac{f'(x)-f'(0)}{2\theta x}= \frac{1}{2}f''(0)\cdot\lim_{x\to0}\frac{1}{\theta},
              \end{aligned}
          \]

          由$f''(0)\neq0$，得$\lim_{x\to0}\theta=\frac{1}{2}$.

          （2）由题设可以写出$\arctan x-0=\frac{x}{1+\xi^{2}}=\frac{x}{1+(\theta x)^{2}}$，所以
          \[
              \theta^{2} = \frac{\frac{x}{\arctan x}-1}{x^{2}}= \frac{x-\arctan x}{x^{2}\arctan
                  x}.
          \]
          因此
          \[
              \lim_{x\to0}\theta^{2}=\lim_{x\to0}\frac{x-\arctan x}{x^{2}\arctan x}=
              \lim_{x\to0}\frac{1-\frac{1}{1+x^{2}}}{3x^{2}}=\frac{1}{3},
          \]
          即$\lim_{x\to0}\theta=\frac{\sqrt{3}}{3}$.
\end{enumerate}
\end{document}




